\section{Related work}
\label{sec:related}

\subsection{Parliamentary monitoring}

Organisations that republish and interpret parliamentary records are an established part of civil society. A 2011 survey counted more than 191 of them, monitoring over 80 national parliaments, with Latin America the most active region~\citep{mandelbaum2011strengthening}. The best known republish debates and votes in searchable form: TheyWorkForYou in the United Kingdom~\citep{mysociety_twfy}, GovTrack in the United States~\citep{govtrack_about}, OpenParliament in Canada~\citep{openparliament_about} and abgeordnetenwatch in Germany~\citep{abgeordnetenwatch_about}. Voteview publishes every roll call of the US Congress with estimated ideal points~\citep{voteview_about}.

Brazil has a rich tradition of its own. Opera\c{c}\~ao Serenata de Amor uses open data to flag suspicious expense reimbursements of legislators~\citep{serenata2016,coutinho2021public}. Radar Parlamentar maps parties by a principal-component analysis of roll calls~\citep{leite2016analise}. The newspaper tool Bas\^ometro measures how often each deputy votes with the government's orientation, with data back to 2003~\citep{estadao_basometro}, and is the closest precedent for the alignment metric we publish. Ranking dos Pol\'iticos~\citep{rankingpoliticos2026criterios} and the \'Indice Legisla Brasil~\citep{legislabrasil2025relatorio} score legislators. Parlametria aggregated legislative data for advocacy organisations from 2019 to 2023~\citep{okbr_parlametria}. Radar do Congresso offers search and analysis of legislators' records~\citep{congressoemfoco_radar}, and Querido Di\'ario republishes municipal official gazettes in open form~\citep{okbr_queridodiario}.

Two observations about this field motivate our work. First, on the public pages we could read of the projects named above (Parlametria's are offline), none describes a way for a reader to verify the published data independently: no digests, signatures, timestamps or transparency logs. The reader is asked to trust the project, its code where it is open, and its hosting. Second, monitoring projects end. VoteWatch Europe stopped in 2022 and its site now carries only a closure notice~\citep{votewatch_closure}; Parlametria's sites are offline; even the domain that hosted the international Declaration on Parliamentary Openness~\citep{openingparliament2012declaration} now serves unrelated content. When such a project disappears, so does the public's ability to check what it once published. Proofs that outlive the project address exactly this.

\subsection{Scores and rankings of legislators}

Rankings are the most popular and the most criticised product of monitoring organisations. Ranking dos Pol\'iticos assigns 75\% of a legislator's score to votes, where a council of the organisation decides which way each vote should count~\citep{rankingpoliticos2026criterios}; critics argue that such criteria embed their authors' world view while presenting themselves as neutral~\citep{queiroz2022rankings,filho2018engodo}. The \'Indice Legisla Brasil combines 17 indicators and in its 2024 edition compresses extreme values so that filing very many bills does not inflate a score~\citep{legislabrasil2025relatorio}.

The experience of TheyWorkForYou is the cautionary tale. Members of Parliament ``increased the number of times they spoke just to increase their scores''~\citep{korthagen2020parliamentary}; the site removed its explicit rankings in 2006 and all comparative statistics in 2020, concluding that ``if we do not believe data can help make meaningful evaluations or requires large qualifications, we should not highlight it''~\citep{parsons2020publishing}. This is Campbell's law~\citep{campbell1979assessing}: an indicator used for social decisions comes under pressure that corrupts both the indicator and the process it measures.

The project's own unpublished score does not escape this critique, although it is a deterministic function of official outcomes with no discretionary input and reports how each rank moves under alternative weights.

\subsection{Roll-call analysis in Brazil}

Political science has studied Brazilian roll calls for three decades. \citet{figueiredo1999executivo} used them to show that parties in the Chamber are disciplined and that the executive controls the agenda, against the view that open-list rules produce only individualists~\citep{ames2001deadlock,mainwaring1999rethinking}. \citet{zucco2011distinguishing} show that voting in the Chamber is structured less by left--right ideology than by a government--opposition dimension, on which parties jump when they enter or leave the coalition. This is the scholarly basis for publishing coincidence with the Government's orientation as the first metric, and the evaluation shows the jumps directly. Our metric is far simpler than ideal-point models~\citep{poole1997congress,clinton2004statistical} or survey-based placements~\citep{power2009estimating}. The simplicity is intended: a ratio of two counts can be recomputed by a reader with a spreadsheet, which a Bayesian item-response model cannot.

\subsection{Cryptographic provenance and verifiable records}

The components we use are well established. Tamper-evident logs~\citep{crosby2009tamper} became public infrastructure with Certificate Transparency~\citep{rfc6962,laurie2014ct}. Sigstore applies the same idea to software signing with short-lived, identity-bound certificates~\citep{newman2022sigstore}; in-toto and SLSA define what a build should attest~\citep{torresarias2019intoto,slsa12}; and The Update Framework showed how to survive key compromise in software distribution~\citep{samuel2010tuf}. Trusted timestamping predates blockchains~\citep{haber1991timestamp,rfc3161}, and Bitcoin has been used to timestamp documents since at least 2015~\citep{gipp2015timestamping}.

These tools were built for software supply chains. Their use for \emph{data} is recent: Sigstore-based model signing covers machine-learning artifacts and datasets by manifest~\citep{openssf-oms}, the SCITT architecture generalises transparency logs to arbitrary signed statements~\citep{rfc9943}, and C2PA attaches signed provenance to media~\citep{c2pa25}. All of them assume that the party vouching for the content is its producer.

The closest practices on the archiving side do not involve the producer either. ``Git scraping'' runs a scheduled CI job that fetches a source and commits the result, so that the repository's history records every change~\citep{willison2020git}; our collector is an instance of it, with attestation and timestamping added. Software Heritage archives public source code and its history under intrinsic, hash-based identifiers~\citep{dicosmo2017software}. Web archives in the WACZ format can carry a signature over the capture~\citep{wacz-signing}, which binds an archive to the party that made it but, like our attestations, cannot show what the original server sent.

Closer to our subject, several systems protect public records. ARCHANGEL uses a distributed ledger run by archives to guarantee the long-term integrity of documents in their custody~\citep{collomosse2018archangel}. Estonia protects government registries with a keyless signature infrastructure based on hash trees~\citep{buldas2013ksi,eestonia-ksi}. \citet{beris2018modeling} prototype the anchoring of Greek government decisions in Bitcoin with a verifier for citizens, and \citet{arruzzoli2026multichain} propose a multi-chain architecture for Italian legislative documents. In each case the custodian of the records, or a system acting for it, applies the protection. A critical survey of blockchain projects in government finds many benefits claimed and few demonstrated~\citep{olnes2017blockchain}.

Our situation is different, and we believe more common: the custodian publishes open data without signatures, and an outsider with no authority republishes it. We found no prior system that combines content-addressed captures, CI-identity attestations in a public transparency log, blockchain timestamps and replayable metrics for public legislative data. Our search covered the websites of the monitoring projects above and literature searches for combinations of transparency logs, Sigstore or blockchain timestamps with legislative, parliamentary or government open data; finding no counter-example is weaker than showing that none exists. We also use no ledger of our own, no token and no smart contract: only hashing, a public signing service and a timestamp.

The closest conceptual relative is end-to-end verifiable voting~\citep{benaloh2015e2e}, whose principle is that anyone can check an election outcome without trusting the software or the officials. Researchers who tested Brazil's voting machines have reported software vulnerabilities in them~\citep{aranha2019return}, and the election disputes described above turned on whether citizens could check the count. Roll-call data are an easier case, since votes on the floor are public and privacy does not constrain the design. The hard part moves elsewhere, to the authenticity of the source.
